\documentclass{article}
\usepackage[T1]{fontenc}
\usepackage{lmodern}
\usepackage{graphicx}
\usepackage{amsmath,amssymb}
\usepackage[a4paper,margin=2cm]{geometry}
\usepackage[absolute,overlay]{textpos}
\setlength{\TPHorizModule}{1mm}
\setlength{\TPVertModule}{1mm}
\usepackage[indonesian]{babel}
\usepackage{hyperref}
\setlength{\emergencystretch}{2em}
% unit: Halaman 48

\begin{document}

% === Halaman 48 (llm) ===

\begin{itemize}
    \item Metode kedua adalah dengan metode Kolbitz. Langkah-langkahnya adalah sebagai berikut:
    \begin{enumerate}
        \item Pilih sebuah kurva eliptik $y^2 \equiv x^3 + ax + b \pmod{p}$ yang mengandung $N$ buah titik.
        \item Misalkan karakter-karakter penyusun pesan adalah angka $0, 1, 2, \dots, 9$ dan huruf $A, B, C, \dots, Z$ yang dikodekan menjadi $10, 11, \dots, 35$.
        \item Kodekan setiap karakter di dalam pesan menjadi nilai $m$ di antara $0$ dan $35$.
        \item Pilih sebuah bilangan bulat $k$ sebagai parameter basis (disepakati kedua pihak).
        \item Untuk setiap nilai $mk$, nyatakan $x = mk + 1$, sulihkan $x$ ke dalam $y^2 \equiv x^3 + ax + b \pmod{p}$ lalu tentukan nilai $y$ yang memenuhi.
        \item Jika tidak ada nilai $y$ yang memenuhi, coba untuk $x = mk + 2, x = mk + 3$, dst, sampai $y^2 \equiv x^3 + ax + b \pmod{p}$ dapat dipecahkan.
        \item Pada proses \textit{decoding}, untuk titik $(x, y)$, tentukan nilai $m$ terbesar tetapi lebih kecil dari $(x - 1)/k$. Kodekan titik $(x, y)$ menjadi symbol $m$.
    \end{enumerate}
\end{itemize}

\end{document}
